%%%%%%%%%%%%%%%%%%%%%%%%%%%%%%%%%%%%%%%%%%%%%%%%%%%%%%%%%%%%%%%%%%%%%%%
%%  Der Einstiegspunkt des Latex-Dokuments                           %%
%%  Diese Datei führt alle Quelldokumente zusammen.                  %%
%%                                                                   %%
%%  Durch den Nutzer durchzuführende Änderungen                      %%
%%  sind entsprechend markiert.                                      %%
%%                                                                   %%
%%%%%%%%%%%%%%%%%%%%%%%%%%%%%%%%%%%%%%%%%%%%%%%%%%%%%%%%%%%%%%%%%%%%%%%
\documentclass[
paper=A4,            %a4paper, alle weiteren Papierformat einstellbar
11pt,                % Schriftgröße (12pt, 11pt (Standard))
BCOR=12mm,           % Bindekorrektur, bspw. 1 cm
headsepline,         % Trennline zum Seitenkopf  
footsepline,         % Trennline zum Seitenfuß
index=totoc,
listof=totoc,
bibliography=totoc,
]
{scrbook}

% 1) pdflatex main
% 2) makeindex.exe %.idx -s etti4.ist
% 3) biber main
% 4) pdflatex main x 2
%%%%%%%%%%%%%%%%%%%%%%%%%%%%%%%%%%%%%%%%%%%%%%%%%%%%%%%%%%%%%%%%%%%%%%%%%%%%%%
\usepackage{lipsum} % generate random text

%-----------------------------------------------------------------------------
% Deutsche Anpassungen
%-----------------------------------------------------------------------------
\usepackage[ngerman]{babel}      % deutsch content->Inhaltsverzeichnis ... 
\RequirePackage[babel]{microtype}
\usepackage[utf8]{inputenc}
%\usepackage[latin1]{inputenc}
\usepackage{inputenc}            % Umlaute
\usepackage{amssymb}             % Pfeile in math 
\usepackage{trfsigns}            % laplace ... Knoten
\RequirePackage{etex}
\usepackage{prettyref}
\usepackage{titleref}


% Font festlegen!
\RequirePackage{courier}								% Font: Schriftart Courier | Setzt Courier als True-Type-Schriftart (ttfamily)
% Change to Helvetica (close to Arial)
\RequirePackage[scaled]{helvet}
\RequirePackage[T1]{fontenc}

\renewcommand\familydefault{\sfdefault}

\usepackage{csquotes}
\usepackage[toc, acronym, nopostdot]{glossaries}
\usepackage{inconsolata}

\RequirePackage[absolute]{textpos}							    % Layout: Absolute Ausrichtung von Text
\RequirePackage[bottom]{footmisc}								% Fußnoten: Positionierung und Verhalten der Fußnoten ändern

\addto\captionsngerman{
  \renewcommand{\figurename}{Abbildung}
  \renewcommand{\tablename}{Tabelle}
  \renewcommand{\lstlistingname}{Quellcode}
}

%-----------------------------------------------------------------------------
% Grafik und Farben
%-----------------------------------------------------------------------------
%\usepackage{showframe} % <--------------------------------------------------- Ausblenden nicht vergessen! 
\usepackage{graphicx}            % Laden von Grafiken
\graphicspath{{./images/}}       % Pfad zu den Bildern
\usepackage{tikz}                % Grafik mit tex erzeugen
\usepackage{multirow}
% \documentclass[xcolor=table]{beamer}


% definierte Farben
% black, blue, brown, cyan, darkgray, gray, green, lightgray, lime, magenta,
% olive, orange, pink, purple, red, teal, violet, white, yellow
% Apricot 	  	  	  	Aquamarine          Bittersweet 	  	  	BlueGreen 
% BlueViolet   	  	    BrickRed            BurntOrange             CadetBlue
% CarnationPink         Cerulean 	  	  	CornflowerBlue          Dandelion
% DarkOrchid 	  	  	Emerald             ForestGreen 	  	  	Fuchsia
% Goldenrod 	  	  	GreenYellow
% JungleGreen 	  	  	Lavender            LimeGreen 	  	  	  	
% Mahogany 	  	  	  	Maroon              Melon 	  	  	  	    MidnightBlue
% Mulberry 	  	  	  	NavyBlue            OliveGreen 	  	  	  	
% OrangeRed 	  	  	Orchid              Peach 	  	  	  	    Periwinkle
% PineGreen 	  	   	Plum                ProcessBlue 	  	  	
% RawSienna 	  	   	RedOrange 	  	  	RedViolet
% Rhodamine 	  	  	RoyalBlue           RoyalPurple 	  	  	RubineRed
% Salmon 	  	  	  	SeaGreen            Sepia 	  	  	  	    SkyBlue
% SpringGreen 	  	  	Tan                 TealBlue 	  	  	  	Thistle
% Turquoise 	  	  	VioletRed 	  	  	  	
% WildStrawberry 	  	YellowGreen 	  	YellowOrange

% RGB define new colors
\definecolor{dkgreen}{rgb}{0,0.6,0}
\definecolor{mauve}{rgb}{0.58,0,0.82}
\definecolor{darkblue}{rgb}{0.1,0,1.0}
\definecolor{darkgreen}{rgb}{0.1,0.8,0.2}
\definecolor{lightblue}{rgb}{0.8,0.85,1}
\definecolor{blue2}{RGB}{0,170,255}
\definecolor{pale turquoise}{rgb}{0.686,0.933,0.933}
\definecolor{light-gray}{gray}{0.95}
\definecolor{ettiOrange}{HTML}{ED6D00}

% Programmcode Farben
\definecolor{mGreen}{rgb}{0,0.6,0}
\definecolor{mGray}{rgb}{0.5,0.5,0.5}
\definecolor{mPurple}{rgb}{0.58,0,0.82}
\definecolor{backgroundColour}{RGB}{240,248,255}

% reuse and mix colors
\colorlet{vlightgray}{lightgray!20}

% mapping - nur die folgenden Farben verwenden
\colorlet{colBackground}{vlightgray}
\colorlet{colPrimary}{blue2}
\colorlet{colSuccess}{green}
\colorlet{colDanger}{red}
\colorlet{colWarning}{magenta!80}
\colorlet{colGrid}{lightgray!40}

% defaults for listings
\colorlet{colKeyword}{magenta}
\colorlet{colComment}{dkgreen}
\colorlet{colString}{mauve}

% defaults for Inhaltsverzeichnis und chapterthumb
\colorlet{colScriptContent}{black}
\colorlet{colScriptCT}{lightblue}



%-----------------------------------------------------------------------------
% Layout
%-----------------------------------------------------------------------------
\parindent0cm
\usepackage{enumitem} % kleinere Abstände

\usepackage{geometry}

			
%\geometry{includehead,includefoot,inner=2.5cm,outer=1.5cm,top=1.5cm,bottom=1.0cm}
\geometry{includehead,includefoot, inner=2.25cm,outer=2.25cm,top=1.5cm,bottom=1.5cm}
\setlength{\footheight}{20.4pt}
\setlength{\headheight}{20.4pt}

\usepackage{siunitx}                 %Mathe Komma einstellen! 
\sisetup{
    detect-family = true,
    %detect-inline-family = math,
    output-decimal-marker = {,}
}


\usepackage{needspace} % Seitenumbruch falls notwendiger Platz nicht mehr vorhanden


\usepackage{rotating}

\usepackage{scrlayer-scrpage}
\pagestyle{scrheadings}
\RequirePackage{xspace}     % Intelligentes Leerzeichen

%-----------------------------------------------------------------------------
% Kopf- und Fußzeile manipulieren
%-----------------------------------------------------------------------------

\renewcommand{\headfont}{\normalfont\sffamily\itshape}          % Kolumnentitel serifenlos
\renewcommand{\pnumfont}{\normalfont\rmfamily}         % Seitennummern typewriter


%-----------------------------------------------------------------------------
% Tabellen
%-----------------------------------------------------------------------------
\usepackage{colortbl}  % Farbige Tabellen
%\usepackage{longtable} % Tabellen die über mehr als eine Seite gehen
\usepackage{booktabs} 
\usepackage{tabularx}
\RequirePackage{tabu}	% Listen:   Erweiterte Listenformatierung tabu Umgebung
\setlength\arrayrulewidth{0.5pt}
\usepackage{makecell}       %Zeilenumbrüche in Tabellen einfügen


% ZUM EINFÜGEN VON BILDERN in TABELLEN
% neuer Befehl: \includegraphicstotab[..]{..}
% Verwendung analog wie \includegraphics
\newlength{\myx} % Variable zum Speichern der Bildbreite
\newlength{\myy} % Variable zum Speichern der Bildhöhe

\newcommand\includegraphicstotab[2][\relax]{%
	% Abspeichern der Bildabmessungen
	\settowidth{\myx}{\includegraphics[{#1}]{#2}}%
	\settoheight{\myy}{\includegraphics[{#1}]{#2}}%
	% das eigentliche Einfügen
	\parbox[c][1.1\myy][c]{\myx}{%
		\includegraphics[{#1}]{#2}}%
}% Ende neuer Befehl


%-----------------------------------------------------------------------------
% Literaturverzeichnis
%-----------------------------------------------------------------------------
\usepackage[
backend=biber,
style=numeric, % style=alphabetic authoryear numeric,
citestyle=numeric,sorting=none,sortcites=true,autopunct=true,
hyperref=true,abbreviate=false,backref=true,block=ragged,
]{biblatex}
\addbibresource{./Literatur.bib} % Literaturdatenbank
\defbibheading{bibempty}{}


%-----------------------------------------------------------------------------
% Index
%-----------------------------------------------------------------------------
\usepackage{makeidx} % Für das Erstellen von einem Index
\makeindex


%-----------------------------------------------------------------------------
% Datum und Zeit
%-----------------------------------------------------------------------------
\usepackage{scrtime}
%-----------------------------------------------------------------------------
% Text
%-----------------------------------------------------------------------------
%\usepackage{microtype} % besserer Wörterabstand und weniger Textbreitenfehler

\RequirePackage[T1]{fontenc}
\RequirePackage{lmodern}
\usepackage{textcomp}
%\usepackage[scale=0.98,ttdefault]{AnonymousPro}
\usepackage{times}

\usepackage{exscale}   % schöne Formeln in PDFs

%-----------------------------------------------------------------------------
% Verlinkungen im Text
%-----------------------------------------------------------------------------
\usepackage[pdftex,bookmarks=true]{hyperref}
\hypersetup{
    pdfborder={0 0 0}, %Es wird kein Rahmen um irgendeine Verlinkung angezeigt, Verlinkungen bleiben aber bestehen
    bookmarksopen=true,
    bookmarksnumbered=true,
    pdftoolbar=true,
    pdfmenubar=true,
    pdffitwindow=true,
    pdfstartview={FitBH},
    pdftitle={Bachelorarbeit},
    pdfauthor={Lukas Diedenhofen},
    pdfsubject={Evaluation von Frameworks zur plattformunabhängigen App-Entwicklung},
    pdfcreator={Lukas Diedenhofen},
    pdfproducer={Lukas Diedenhofen},
    pdfkeywords={Dart} {Flutter} {ReactNative} {iOS} {Android} {Bachelorarbeit},
    colorlinks=false,
    linkcolor=red,
    citecolor=green,
    citebordercolor={0 1 0},
    filecolor=magenta,
    urlbordercolor={0 1 1},
    linktoc=all     %Verlinkungen im Inhaltsverzeichnis aktivieren oder speziell einstellen
}
\urlstyle{sf}

%pdftitle={Bachelorarbeit},pdfauthor={Lukas Diedenhofen},pdfcreator={Lukas Diedenhofen},pdfsubject={Evaluation von Frameworks zur plattformunabhängigen App-Entwicklung},

%-----------------------------------------------------------------------------
% Boxen auch zweispaltig
%-----------------------------------------------------------------------------
\usepackage{multicol}
\usepackage[most]{tcolorbox} % all but minted and documentation
%\usepackage[most,skins,listingsutf8]{tcolorbox}

% an dieser Stelle nach tcolorbox !!
\RequirePackage{setspace}
%\usepackages[onehalfspace]{setspace}
%%% [setspace]  %%%%%%%%%%%%%%%%%%%%%%%%%%%%%%%%%%%%%%%%%%%%%%
\onehalfspacing                                % Leerzeilen eineinhalbzeilig



\usepackage{environ}

% standard box
\newtcolorbox{tcbDefault}[1][]{
	colback=blue!5,
	colframe=blue!50,
	#1
}

% standard box with title
\newtcolorbox{tcbDefaultT}[2][]{
	coltitle=blue!50!black,
	colback=blue!5,
	colframe=blue!35,
	title={#2},
	%fonttitle=\bfseries,
	#1
}

%-----------------------------------------------------------------------------
% Listings with tcb
%-----------------------------------------------------------------------------

\newtcolorbox[auto counter,number within=section,]{tcbListingDefault}[2][]{
	enhanced,
	sharp corners=downhill,
	bicolor,
	arc=.5cm,
	top       = \tcboxedtitleheight,
	bottom    = 5mm,
	attach boxed title to top right={yshift=-\tcboxedtitleheight},
	boxed title style={
		size=small,
		colback=gray!20,
		colframe=gray!30,
		sharp corners=downhill,
		arc=.5cm,
		top=1mm,bottom=1mm,left=1mm,right=1mm
	},
	colframe  = gray!30,
	colback   = white,
	colbacklower=gray!5,
	coltitle = blue!50!black,
	fonttitle=\bfseries,
	leftupper=6mm,
	rightupper=0mm,
	leftlower=0mm,
	%boxsep=5mm, 
	middle=5mm,
	title=Listing~\thetcbcounter~ #2,
	#1
}

% with option side by side
\newtcolorbox[auto counter,number within=section,]{tcbListingDefaultSBS}[3][]{
	enhanced,
	sharp corners=downhill,
	bicolor,
	sidebyside,
	arc=.5cm,
	top       = \tcboxedtitleheight,
	bottom    = 0mm,
	attach boxed title to top right={yshift=-\tcboxedtitleheight},
	boxed title style={
		size=small,
		colback=colBackground,
		colframe=gray!30,
		sharp corners=downhill,
		arc=.5cm,
		top=1mm,bottom=1mm,left=1mm,right=1mm
	},
	colframe  = gray!30,
	colback   = white,
	colbacklower=gray!5,
	coltitle = blue!50!black,
	fonttitle=\bfseries,
	righthand width=#3,
	leftupper=6mm,
	rightupper=0mm,
	leftlower=0mm,
	boxsep=0mm, 
	middle=0mm,
	title=Listing~\thetcbcounter~ #2,
	#1
}


%\newtcbinputlisting[use counter=tcbListing]{\tcbinpListing}[3][]{%
\newtcbinputlisting[use counter from=tcbListingDefault]{\tcbinpListing}[3][]{%
	%\newtcbinputlisting[auto counter,number within=section]{\tcbinpListing}[3][]{%
	enhanced,
	bicolor,
	attach boxed title to top right={yshift=-\tcboxedtitleheight},
	boxed title style={
		size=small,
		colback=gray!20,
		colframe=gray!30,
		sharp corners=downhill,
		arc=.5cm,
		top=1mm,bottom=1mm,left=1mm,right=1mm
	},
	listing options={numbers=left,numberstyle=\tiny\color{red!75!black}},
	arc=.5cm,
	%left=6mm,      % nr ausserhalb Box
	boxsep=3mm,    % nr im Rand
	%fonttitle=\color{black}\itshape\ttfamily,
	fonttitle = \color{black}\bfseries,
	colframe  = gray!30,
	colback   = white,%gray!20,
	top       = \tcboxedtitleheight,
	bottom    = 0mm,
	sharp corners=downhill,
	listing file = {#2},
	%title        = Listing (\thetcbcounter) of \texttt{#2} #3,
	title        = Listing \thetcbcounter: #3,
	listing only,
	colbacklower=gray!5,
	breakable,
	#1
}

% Farben für die Programmierung festlegen:
\definecolor{VC_Blue}{RGB}{64,41,255}
\definecolor{VC_Green}{RGB}{42,139,15}
\definecolor{VC_BlueGrey}{RGB}{60,120,147}
\definecolor{VC_Brown}{RGB}{148,116,54}
\definecolor{VC_Blue2}{RGB}{12,113,197}
\definecolor{VC_Red}{RGB}{254,25,0}
\definecolor{VC_Green2}{RGB}{37,130,0}
\definecolor{VC_DarkBlue}{RGB}{12,24,128}
\definecolor{VC_Purpel}{RGB}{177,0,218}
\definecolor{VC_BrownRed}{RGB}{170,40,20}
\definecolor{VC_Dart_main}{RGB}{131,94,38}
\definecolor{UNIBW2}{RGB}{248,115,4}
\definecolor{UNIBW}{RGB}{236,111,0}
\definecolor{NumbersDart}{RGB}{15,130,101}
\colorlet{headerUNI}{UNIBW!60}
\colorlet{dunkelUNI}{gray!20}
%\colorlet{dunkelUNI}{UNIBW!25}
\colorlet{hellUNI}{white}
%\colorlet{hellUNI}{UNIBW!7}
\colorlet{dunkelZeile}{gray!20}


%-----------------------------------------------------------------------------
% Listings und pdf include
%-----------------------------------------------------------------------------
\usepackage{pdfpages}       % include pdf's 
\usepackage{listingsutf8}   % utf-8 encoded source files
\usepackage{arydshln}


% Quellcode: Eigenes Highlightening und Formatieren des Sourcecodes:
\RequirePackage{listings}

%-----------------------------------------------------------------------------
% Style for JavaScript und ReactNative
%-----------------------------------------------------------------------------

\definecolor{RN_Numbers}{RGB}{8,134,89}
\definecolor{RN_Key}{RGB}{0,0,255}
\definecolor{RN_Function}{RGB}{126,100,40}
\definecolor{VC_Brown}{RGB}{148,116,54}
\definecolor{VC_Blue2}{RGB}{12,113,197}
\definecolor{VC_Red}{RGB}{254,25,0}
\definecolor{VC_Green2}{RGB}{37,130,0}
\definecolor{VC_DarkBlue}{RGB}{12,24,128}
\definecolor{VC_Purpel}{RGB}{177,0,218}
\definecolor{VC_BrownRed}{RGB}{170,40,20}
\definecolor{VC_Dart_main}{RGB}{131,94,38}
\definecolor{UNIBW2}{RGB}{248,115,4}
\definecolor{UNIBW}{RGB}{236,111,0}
\definecolor{NumbersDart}{RGB}{15,130,101}

\lstdefinestyle{Java}{
    literate={0}{{\textcolor{RN_Numbers}{0}}}{1}%
                 {1}{{\textcolor{RN_Numbers}{1}}}{1}%
                 {2}{{\textcolor{RN_Numbers}{2}}}{1}%
                 {3}{{\textcolor{RN_Numbers}{3}}}{1}%
                 {4}{{\textcolor{RN_Numbers}{4}}}{1}%
                 {5}{{\textcolor{RN_Numbers}{5}}}{1}%
                 {6}{{\textcolor{RN_Numbers}{6}}}{1}%
                 {7}{{\textcolor{RN_Numbers}{7}}}{1}%
                 {8}{{\textcolor{RN_Numbers}{8}}}{1}%
                 {9}{{\textcolor{RN_Numbers}{9}}}{1}%
                 {.0}{{\textcolor{RN_Numbers}{.0}}}{2}% Following is to ensure that only periods
                 {.1}{{\textcolor{RN_Numbers}{.1}}}{2}% followed by a digit are changed.
                 {.2}{{\textcolor{RN_Numbers}{.2}}}{2}%
                 {.3}{{\textcolor{RN_Numbers}{.3}}}{2}%
                 {.4}{{\textcolor{RN_Numbers}{.4}}}{2}%
                 {.5}{{\textcolor{RN_Numbers}{.5}}}{2}%
                 {.6}{{\textcolor{RN_Numbers}{.6}}}{2}%
                 {.7}{{\textcolor{RN_Numbers}{.7}}}{2}%
                 {.8}{{\textcolor{RN_Numbers}{.8}}}{2}%
                 {.9}{{\textcolor{RN_Numbers}{.9}}}{2}%
                 ,
    language        =   Java,
    keywordstyle    =   \color{VC_Blue},
    commentstyle    =   \color{VC_Green},
    stringstyle     =   \itshape\color{VC_BrownRed},
    backgroundcolor =   \color{white},
    numberstyle     =   \small\ttfamily\color{VC_BlueGrey},
    basicstyle      =   \small\ttfamily\color{black},
    %basicstyle      =   \small\ttfamily\color{VC_DarkBlue},
    xleftmargin     =   25pt,
    emph            =   [1]{import, from, as, return, if, else, export, default, await},
    emphstyle       =   [1]\color{VC_Purpel},
    emph            =   [2]{*, const, function, false, null, =>, let, class, extends, async, true, this, super, constructor },
    emphstyle       =   [2]\color{VC_Blue},
    emph            =   [3]{[0-9]},
    emphstyle       =   [3]\color{VC_Red},
    emph            =   [4]{HelloWorldApp, View, Icon, Text, Image, ScrollView, Firebase, TouchableOpacity, SafeAreaView, StatusBar, StyleSheet, TextInput, Button, Alert},
    emphstyle       =   [4]\color{VC_BlueGrey},
    emph            =   [5]{style, name, size, source, contentContainerStyle, component, options, visible, onPress, disabled, author, barStyle, placeholder, keyboardType, underlineColorAndroid, onChangeText, title },
    emphstyle       =   [5]\color{VC_Red},
    emph            =   [6]{Platform, styles, },
    emphstyle       =   [6]\color{VC_Blue2},
    emph            =   [7]{count, state, props, buttonPlusFive, buttonPlusTwo, buttonPlusOne, buttonMinusFive, buttonMinusTwo, buttonMinusOne, textPlusOne, textPlusTwo, textPlusFive, textMinusOne, textMinusTwo, countContainer, text, counter, fixToText, container, OS, flex, justifyContent, paddingHorizontal, alignItems, backgroundColor, padding, width, height, shadowColor, shadowRadius, shadowOpacity, shadowOpacity, shadowOffset, elevation, color, fontSize, flexDirection, textReset, resetView, buttonReset, linksrechts, date, amount, category},
    emphstyle       =   [7]\color{VC_DarkBlue},
    emph            =   [8]{increment, setState, decrement, create, render, decrementTwo, decrementFive, incrementTwo, incrementFive, reset, setAddItemScreen, alert, addItem, onChangeAmount },
    emphstyle       =   [8]\color{RN_Function},
    frame           =   l,
}


%-----------------------------------------------------------------------------
% Style for Dart & flutter
%-----------------------------------------------------------------------------

\lstdefinestyle{Dart}{
    literate={0}{{\textcolor{NumbersDart}{0}}}{1}%
             {1}{{\textcolor{NumbersDart}{1}}}{1}%
             {2}{{\textcolor{NumbersDart}{2}}}{1}%
             {3}{{\textcolor{NumbersDart}{3}}}{1}%
             {4}{{\textcolor{NumbersDart}{4}}}{1}%
             {5}{{\textcolor{NumbersDart}{5}}}{1}%
             {6}{{\textcolor{NumbersDart}{6}}}{1}%
             {7}{{\textcolor{NumbersDart}{7}}}{1}%
             {8}{{\textcolor{NumbersDart}{8}}}{1}%
             {9}{{\textcolor{NumbersDart}{9}}}{1}%
             {.0}{{\textcolor{NumbersDart}{.0}}}{2}% Following is to ensure that only periods
             {.1}{{\textcolor{NumbersDart}{.1}}}{2}% followed by a digit are changed.
             {.2}{{\textcolor{NumbersDart}{.2}}}{2}%
             {.3}{{\textcolor{NumbersDart}{.3}}}{2}%
             {.4}{{\textcolor{NumbersDart}{.4}}}{2}%
             {.5}{{\textcolor{NumbersDart}{.5}}}{2}%
             {.6}{{\textcolor{NumbersDart}{.6}}}{2}%
             {.7}{{\textcolor{NumbersDart}{.7}}}{2}%
             {.8}{{\textcolor{NumbersDart}{.8}}}{2}%
             {.9}{{\textcolor{NumbersDart}{.9}}}{2}%
             ,
    language        =   c,
    keywordstyle    =   \color{VC_Blue},
    commentstyle    =   \color{VC_Green},
    stringstyle     =   \itshape\color{VC_BrownRed},
    backgroundcolor =   \color{white},
    numberstyle     =   \small\ttfamily\color{VC_BlueGrey},
    basicstyle      =   \small\ttfamily\color{black},
    xleftmargin     =   25pt,
    emph            =   [1]{return, if, for},
    emphstyle       =   [1]\color{VC_Purpel},
    emph            =   [2]{import, void, class, extends, @override, final, with, false, super, this, get, null, var, @required},
    emphstyle       =   [2]\color{VC_Blue},
    emph            =   [3]{StatefulWidget, Widget, Function, NewTransaction, MyApp, BuildContext, MaterialApp, Colors, ThemeData, AppBar, List,                                    Transaction, bool, double, String, Text, StatelessWidget, SafeArea, Column, Scaffold, Platform, Alignment, Radius, Container, TextStyle, TextAlign, EdgeInserts, RoundedRectangleBorder, BorderRadius, EdgeInsets, Ink, BoxDecoration, LinearGradient, Theme, BoxConstraints, RaisedButton, CupertinoPageScaffold, FloatingActionButton, Icon, Icons, FloatingActionButtonLocation, CupertinoButton, FontWeight, ElevatedButton},
    emphstyle       =   [3]\color{VC_BlueGrey},
    emph            =   [4]{main, build, runApp, light, copyWith, createState, initState, addObserver, print, dispose, toList, now, isAfter, where,                                 subtract, add, setState, _incrementCounter, of, all, elliptical, _startAddNewTransaction, _submitData, _excrementCounterFive},
    emphstyle       =   [4]\color{VC_Dart_main},
    emph            =   [5]{1,2,3,4,5,6,7,8,9,0},
    emphstyle       =   [5]\color{UNIBW},
    frame           =   l,
}


% Style für C
\lstdefinestyle{c}{
    commentstyle=\color{mGreen},
    keywordstyle=\color{magenta},
    stringstyle=\color{mPurple},
    breakatwhitespace=false,         
    breaklines=true,                 
    captionpos=b,                    
    keepspaces=false,                 
    numbers=left,                    
    numbersep=5pt,                  
    showspaces=false,                
    showstringspaces=false,
    showtabs=false,                  
    tabsize=2,
    language=C,
    emphstyle={\bfseries\color{red}},
}


% Einstellung des Default-Styles:
\lstset{ 
    basicstyle  = \small,              % Die Textgröße für den Quelltext
	numbers       = left,              % Platzierung Zeilennummern
	numberstyle   = \tiny\color{gray}, % Stil für die Seitennummern
	stepnumber    = 1,                 % Schritt zwischen nummerierten Zeilen.
	backgroundcolor= \color{white},    % Hintergrundfarbe für den Quelltext
	showspaces    = false,             % show spaces adding particular underscores
	showstringspaces = false,          % underline spaces within strings
	showtabs      = false,             % show tabs within strings adding particular underscores
	frame         = none,              % adds a frame around the code
	% lines, single, none
	rulecolor     = \color{black},     % if not set, the frame-color may be changed on
	% line-breaks within not-black text (e.g. comments (green here))
	tabsize       = 2,                 % Die Größe für einen Tabulator
	captionpos    = b,                 % Die Position der Überschrift. Hier b = Bottom
	breaklines    = true,              % Automatischer Zeilenumbruch
	breakatwhitespace=false,           % sets if automatic breaks only at whitespace
	title=\lstname,                    % show the filename of files 
	% WEGLASSEN damit lstlistings ohne caption weniger Platz brauchen
	keywordstyle  = \color{colKeyword},      % keyword style
	commentstyle  = \color{colComment}\ttfamily,   % comment style
	stringstyle   = \color{colString},     % string literal style
	escapeinside  = {\%*}{*)},         % if you want to add LaTeX within your code
	%linerange    = {1-13},            % Zeilenauswahl
	%linewidth    = 16cm               % = 0.9\linewidth, fix
	%xleftmargin  = 1.5cm,             % Einrücken von rechts
	%xrightmargin = 3.5cm,             % Einrücken von links
	%left         = 15pt,               % nr ausserhalb Box
	%boxsep       = 0pt,               % nr im Rand
	%boxrule      = 0pt,               % ohne Rand
	prebreak = \raisebox{0ex}[0ex][0ex]{\ensuremath{\hookleftarrow}},
	%	inputpath     = ./Source	         % default source input path
	% Achtung nicht inputpath verwenden sonst Konflikt mit tcblistings
}
\lstset{literate=
  {á}{{\'a}}1 {é}{{\'e}}1 {í}{{\'i}}1 {ó}{{\'o}}1 {ú}{{\'u}}1
  {Á}{{\'A}}1 {É}{{\'E}}1 {Í}{{\'I}}1 {Ó}{{\'O}}1 {Ú}{{\'U}}1
  {à}{{\`a}}1 {è}{{\`e}}1 {ì}{{\`i}}1 {ò}{{\`o}}1 {ù}{{\`u}}1
  {À}{{\`A}}1 {È}{{\'E}}1 {Ì}{{\`I}}1 {Ò}{{\`O}}1 {Ù}{{\`U}}1
  {ä}{{\"a}}1 {ë}{{\"e}}1 {ï}{{\"i}}1 {ö}{{\"o}}1 {ü}{{\"u}}1
  {Ä}{{\"A}}1 {Ë}{{\"E}}1 {Ï}{{\"I}}1 {Ö}{{\"O}}1 {Ü}{{\"U}}1
  {â}{{\^a}}1 {ê}{{\^e}}1 {î}{{\^i}}1 {ô}{{\^o}}1 {û}{{\^u}}1
  {Â}{{\^A}}1 {Ê}{{\^E}}1 {Î}{{\^I}}1 {Ô}{{\^O}}1 {Û}{{\^U}}1
  {ã}{{\~a}}1 {ẽ}{{\~e}}1 {ĩ}{{\~i}}1 {õ}{{\~o}}1 {ũ}{{\~u}}1
  {Ã}{{\~A}}1 {Ẽ}{{\~E}}1 {Ĩ}{{\~I}}1 {Õ}{{\~O}}1 {Ũ}{{\~U}}1
  {œ}{{\oe}}1 {Œ}{{\OE}}1 {æ}{{\ae}}1 {Æ}{{\AE}}1 {ß}{{\ss}}1
  {ű}{{\H{u}}}1 {Ű}{{\H{U}}}1 {ő}{{\H{o}}}1 {Ő}{{\H{O}}}1
  {ç}{{\c c}}1 {Ç}{{\c C}}1 {ø}{{\o}}1 {å}{{\r a}}1 {Å}{{\r A}}1
  {€}{{\euro}}1 {£}{{\pounds}}1 {«}{{\guillemotleft}}1
  {»}{{\guillemotright}}1 {ñ}{{\~n}}1 {Ñ}{{\~N}}1 {¿}{{?`}}1 {¡}{{!`}}1 
}


%----------------------------------------------------------------------------------------
% Inhaltsverzeichnis 
%----------------------------------------------------------------------------------------

%\usepackage{titletoc} % Manipulieren der Headings
%\contentsmargin{0cm} % Removes the default margin

% part style
%\titlecontents{part}
%[0cm]                       % Left indentation
%{\addvspace{20pt}\bfseries} % Spacing and font options
%{}
%{}
%{}

% Chapter Style
%\titlecontents{chapter}
%[1.25cm]                                                                        % Left indentation
%{\addvspace{10pt}\large\sffamily\bfseries}                                      % Spacing and font options
%{\color{black}\contentslabel[\large\thecontentslabel]{1.00cm}\color{black}}     % Formatting of numbered sections of this type
%{}                                                                              % Formatting of numberless sections of this type
%{\color{black}$\;$\small\titlerule*[.3cm]{ }$\;$\large\thecontentspage}         % filler and page number
%{\color{black}$\;$\small\titlerule*[.3cm]{.}$\;$\large\thecontentspage}

% Section Style
%\titlecontents{section}
%[2.5cm]                                                                         % Left indentation
%{\addvspace{3pt}\normalsize\sffamily}      %\bfseries                            % Spacing and font options
%{\contentslabel[\thecontentslabel]{1.25cm}}                                     % Formatting of numbered sections 
%{}                                                                              % Formatting of numberless sections of this type
%{\color{black}\small\titlerule*[.4cm]{.}$\;$\normalsize\thecontentspage}        % Formatting of the filler and the page number

% Subsection Style
%\titlecontents{subsection}
%[3.75cm]                                                                        % Left indentation
%{\addvspace{1pt}\sffamily\small}                                                % Spacing and font options
%{\contentslabel[\thecontentslabel]{1.25cm}}                                     % Formatting of numbered 
%{}                                                                              % Formatting of numberless sections of this type
%{\color{black}\small\titlerule*[.4cm]{.}$\;$\thecontentspage}                   % Formatting of the filler  and the page number


% · <-- Mittelpunkt
% – <-- Gedankenstrich


\usepackage{pgf-pie}
\usepackage{tikz}
\usepackage{pgfplots}
\pgfplotsset{width=15cm, height=8cm, compat=1.17}

\usepackage{bchart}
\usepackage{caption}          % Laden aller notwendigen Pakte und Definitionen	
\usepackage{tabu} 


%% Packages für Grafiken & Abbildungen 
\usepackage{subfigure} %%Teilabbildungen in einer Abbildung
\usepackage{caption}

\newenvironment{abstract}{}{}
\usepackage{abstract}


\RequirePackage{float}							% Gleitobjekte(table,figure): Erweitertung
\RequirePackage{placeins}						% Gleitobjekte(table,figure): Position einschränken

%% Sonstige Pakete 
\usepackage{eurosym}                            % Euro Symbol
\usepackage{dsfont}                             % Fonts für Zahlenbereiche


%% Abkürzungen
%%%%%%%%%%%%%%%%%%%%%%%%%%%%%%%%%%%%%%%%
%%Acronyms%%%%%%%%%%%%%%%%%%%%%%%%%%%%%%
%%%%%%%%%%%%%%%%%%%%%%%%%%%%%%%%%%%%%%%%
\newacronym{ba}{BA}{Bachelorarbeit}
\newacronym{etti}{ETTI}{Elektrotechnik und technische Informatik}


%%%%%%%%%%%%%%%%%%%%%%%%%%%%%%%%%%%%%%%%
%%Definitions%%%%%%%%%%%%%%%%%%%%%%%%%%%
%%%%%%%%%%%%%%%%%%%%%%%%%%%%%%%%%%%%%%%%
\newglossaryentry{latex}
{
        name=latex,
        description={Is a mark up language specially suited for 
scientific documents}
}
\makeglossaries


\lstdefinelanguage{NOSTYLE}{
  alsodigit = {-},
  keywords = {}
}

\setlength{\parindent}{0em}

\usepackage{pgf-pie}
\usepackage{pgfgantt}
\usepackage{blindtext}
\usepackage{verbatim}



%%%%%%%%%%%%%%%%%%%%%%%%%%%%%%%%%%%%%%%%%%%%%%%%%%%%%%%%%%%%%%%%%%%%%%%%%%%
%% Die folgenden Variablen sind durch den Nutzer anzupassen.             %%
%% Sie werden in den entsprechenden Dokumenten automatisch übernommen;   %%
%% Änderungen in den Dokumenten sind dann nicht mehr nötig.              %%
%%%%%%%%%%%%%%%%%%%%%%%%%%%%%%%%%%%%%%%%%%%%%%%%%%%%%%%%%%%%%%%%%%%%%%%%%%%
\def\theCourseOfStudies{Wehrtechnik - Informations- und Elektrotechnik}
\def\theCertificate{Bachelorarbeit}
\def\theTitle{Einsatz von LLM Tools im Software Engineering Prozess}
\def\theSubTitle{}
\def\theAuthor{Jonas Gundlach}
\def\theAdviserA{Prof. Dr. rer. nat. Andrea Baumann}
\def\theAdviserB{~}
\def\thePublicationDate{XX.XX.XXXX}
\def\thePublicationMY{Juli 2024}
%%%%%%%%%%%%%%%%%%%%%%%%%%%%%%%%%%%%%%%%%%%%%%%%%%%%%%%%%%%%%%%%%%%%%%%%%%%
%%%%%%%%%%%%%%%%%%%%%%%%%%%%%%%%%%%%%%%%%%%%%%%%%%%%%%%%%%%%%%%%%%%%%%%%%%%

\DeclareMathOperator{\ceil}{ceil}
\DeclareMathOperator{\floor}{floor}

\newcounter{savepage}
\begin{document}


%Caption nach unten setzen
\captionsetup[figure]{skip=10pt}
%\captionsetup[table]{skip=10pt}
\captionsetup[lstlisting]{skip=10pt}

\pagenumbering{Roman}

%% Titelseite
%%%%%%%%%%%%%%%%%%%%%%%%%%%%%%%%%%%%%%%%%%%%%%%%%%%%%%%%%%%%%%%%%%%%%%%%%%%
%% Die Titelseite des Dokuments.                                         %%
%% In diesem Dokument sind keine Änderungen durch den Nutzer             %%
%% notwendig; Die Variablen werden aus der thesis.tex übernommen.        %%
%%%%%%%%%%%%%%%%%%%%%%%%%%%%%%%%%%%%%%%%%%%%%%%%%%%%%%%%%%%%%%%%%%%%%%%%%%%
\begin{titlepage}

\setlength{\topmargin}{0pt}

\begin{center}

\begin{center}
  \includegraphics[width=12cm]{images/logos/UniBwMSignet.pdf}
\end{center}

\vspace{30pt}


%% Hier wird der Titel der Arbeit eingetragen
\huge
\textsc{\theTitle}\\
\vspace{5pt}
\large
\textsc{\theSubTitle}\\
\vspace{40pt}

\normalsize
\textsc{\theCertificate}                                % BACHELORARBEIT 
\textsc{\\zur Erlangung des akademischen Grades\\}
\textsc{Bachelor of Engineering (B.Eng.)}\\

\vspace{40pt}


\textrm{\large \theAuthor}\\                            % Lukas Diedenhofen

\normalsize

\vspace{40pt}

Betreuerin:\\

\vspace{0pt}

\textrm{\large \theAdviserA}\\                          % Prof. Dr. rer. nat. Andrea Baumann

\vspace{30pt}

Tag der Abgabe: \thePublicationDate\\ 

\vspace{30pt}

eingereicht bei\\
Universität der Bundeswehr München\\
Fakultät für Elektrotechnik und Technische Informatik\\
\vspace{15pt}
\includegraphics[width=242pt]{images/logos/UniBwM_ETTI_6_geschnitten.pdf}


\vspace{25pt}
Neubiberg, \thePublicationMY 
\end{center}
\end{titlepage}
\cleardoublepage
%%%%%%%%%%%%%%%%%%%%%%%%%%%%%%%%%%%%%%%%%%%%%%%%%%%%%%%%%%%%%%%%%%%%%%%%%%%
%% Die Erklärung de Arbeit.                                              %%
%% In diesem Dokument sind keine Änderungen durch den Nutzer             %%
%% notwendig; Die Variablen werden aus der thesis.tex übernommen.        %%
%%%%%%%%%%%%%%%%%%%%%%%%%%%%%%%%%%%%%%%%%%%%%%%%%%%%%%%%%%%%%%%%%%%%%%%%%%%
\newpage
\chapter*{Erklärung}

\addcontentsline{toc}{chapter}{Erklärung}

\vspace{1cm}

Hiermit versichere ich, dass ich die vorliegende Arbeit selbstständig verfasst, noch nicht anderweitig für Prüfungszwecke vorgelegt und keine anderen als die angegebenen Quellen und Hilfsmittel benutzt habe, insbesondere keine anderen als die angegebenen Informationen. \\
\hspace{\fill}
\vspace{0.5cm}

Der Speicherung meiner Bachelorarbeit zum Zweck der Plagiatsprüfung stimme ich zu. Ich versichere, dass die elektronische Version mit der gedruckten Version inhaltlich übereinstimmt.\\
\hspace{\fill}
\vspace{0.6cm}

Neubiberg, den 23.09.2024 
\hspace{\fill}
\newline\hspace{\fill}\newline\hspace{\fill}\newline\hspace{\fill}\newline\hspace{\fill}
\rule{0.3\textwidth}{0.4pt} \newline
\theAuthor
\hspace{\fill}
\cleardoublepage

%%%%%%%%%%%%%%%%%%%%%%%%%%%%%%%%%%%%%%%%%%%%%%%%%%%%%%%%%%%%%%%%%%%%%%%%%%%
%% Die folgenden Variable ist durch den Nutzer anzupassen.               %%
%% Sollte kein Abstract gewünscht sein, kann die folgende Variable       %%
%% auskommentiert werden.                                                %%
%%%%%%%%%%%%%%%%%%%%%%%%%%%%%%%%%%%%%%%%%%%%%%%%%%%%%%%%%%%%%%%%%%%%%%%%%%%

\begin{abstract}


    %% Anfang des Abstract-Texts
     
    Im Rahmen dieser Arbeit wird untersucht, inwieweit die drei LLM-Tools ChatGPT, Google Gemini und Mistral Le Chat im Bereich des Software 
    Engineerings eingesetzt werden können. Dabei wird betrachtet, wie geeignet die drei unterschiedlichen LLM Tools sind, die in einem Software-Engineering-Prozess 
    verwendeten Dokumente selbstständig zu erstellen. Die Dokumente werden am Beispiel eines „Mensch ärgere dich nicht“-Spiels erstellt.

    Nach einer kurzen Beschreibung der drei LLM-Tools und einem Überblick über die Inhalte der verschiedenen Dokumente wird untersucht, inwieweit 
    die Tools im Rahmen eines Software-Engineering-Prozesses unterstützen können. Unter zahlreichen Unterstützungsmöglichkeiten, wird sich innerhalb der 
    vorliegenden Arbeit auf eine generische Dokumentenerstellung konzentriert.
    
    In der Analyse werden die Ergebnisse der Tools, die auf dem beiliegenden USB-Stick zu finden sind, besprochen und auf ihre Fehler untersucht. 
    Dabei stellen insbesondere das Format der UML-Diagramme und die Konsistenz zwischen den einzelnen Dokumenten eine Herausforderung dar. 
    Anschließend werden auch die Funktionen für die Schnittstelle mit ChatGPT erstellt, da dieses Tool die Dokumente am erfolgreichsten erstellt hat.
    
    Abschließend gibt es eine Zusammenfassung der Ergebnisse und eine Beschreibung, wie gut mit den einzelnen Tools gearbeitet werden konnte. Danach 
    folgt ein Abschnitt zur Kritik an dieser Arbeit und der Arbeitsweise. Besonders aufgrund der vielen unterschiedlichen Aspekte, die bei der 
    Erstellung der Dokumente berücksichtigt werden müssen, und der unterschiedlichen Ausgaben der drei Tools war es schwierig, den Überblick zu 
    behalten. Zum Schluss folgt ein kurzer Ausblick auf offene Punkte, die in diesem Bereich noch untersucht werden können. Da die drei Tools noch 
    relativ neu sind und die Entwicklung in den letzten Jahren erst richtig Fahrt aufgenommen hat, ist zu erwarten, dass sich in diesem Bereich noch 
    viel tun wird.

    %% Ende des Abstract-Texts
    
    
\end{abstract}

%%%%%%%%%%%%%%%%%%%%%%%%%%%%%%%%%%%%%%%%%%%%%%%%%%%%%%%%%%%%%%%%%%%%%%%%%%%
%%%%%%%%%%%%%%%%%%%%%%%%%%%%%%%%%%%%%%%%%%%%%%%%%%%%%%%%%%%%%%%%%%%%%%%%%%%

%% Inhaltsverzeichnis
\setcounter{tocdepth}{2}                %% Festlegen der Tiefe des Inhaltsverzeichnisses
\setcounter{secnumdepth}{2}

\tableofcontents                        %% Erzeugen des Inhaltsverzeichnisses
\cleardoublepage
\pagenumbering{arabic}

%%%%%%%%%%%%%%%%%%%%%%%%%%%%%%%%%%%%%%%%%%%%%%%%%%%%%%%%%%%%%%%%%%%%%%%%%%%
%% Die folgenden Variablen sind durch den Nutzer anzupassen.             %%
%% Hier werden die einzelnen Quelldokumente aus dem sourece-Folder       %%
%% referenziert, um im pdf-Dokument gerendert zu werden.                 %%
%%%%%%%%%%%%%%%%%%%%%%%%%%%%%%%%%%%%%%%%%%%%%%%%%%%%%%%%%%%%%%%%%%%%%%%%%%%
%% Hauptteil der Arbeit 

\chapter{Einleitung} \index{Einleitung}

Die rasante Entwicklung im Bereich der Künstlichen Intelligenz hat in den letzten Jahren bedeutende Fortschritte ermöglicht, 
insbesondere durch den Einsatz von Large Language Models (LLM). Diese mächtigen Tools, die auf umfangreichen Textkorpora 
trainiert wurden, haben das Potenzial, vielfältige Aufgaben im Software Engineering Prozess zu optimieren. 
In dieser Bachelorarbeit wird untersucht, wie LLM-Tools im Software Engineering eingesetzt werden können, um die Effizienz und 
Qualität der Entwicklungsprozesse zu verbessern.

In diesem ersten Kapitel wird zunächst auf die Motivation hinter dieser Arbeit eingegangen. Anschließend wird die Zielsetzung 
und dann die Struktur der Arbeit beschrieben.

\section{Motivation} \index{Motivation} \label{Motivation}

KIs ersetzen irgendwann die Menschen! Dies ist eine Aussage, die man mittlerweile recht häufig hört. Unter anderem durch Filme 
und Serien ist das Ansehen von KIs nicht immer positiv, und einige Menschen sehen eine große Gefahr darin. Besonders durch die 
rasante Entwicklung in den letzten zwei Jahren haben sogenannte Large Language Models (LLMs) viel Aufmerksamkeit auf sich gezogen 
und die Möglichkeiten, zum Beispiel in der Informationsbeschaffung, stark vergrößert. Gerade in Schulen und Universitäten haben 
solche Tools mittlerweile eine große Relevanz erlangt, was nicht von jedem positiv gesehen wird. Viele Pädagogen behaupten, dass 
LLM-Tools schlecht für die Entwicklung von Schülern und Studenten sind, da man nun Hausaufgaben oder ähnliche Arbeiten einfach an 
LLM-Tools abgeben und in wenigen Sekunden automatisch erstellen lassen kann. Doch wie gut sind eigentlich die Ausgaben von LLM Tools?

In einer vorherigen Projektarbeit zum Thema ``Erzeugung von Algorithmen mit LLM-Tools für Java'' wurde bereits untersucht 
wie gut drei unterschiedliche Tools einfache Algorithmen zum Sortieren von Listen und einen etwas komplexeren Algorithmus 
für Spiele wie Tic-Tac-Toe oder Vier Gewinnt erstellen können. Dies hat, zumindest bei den einfacheren Algorithmen, bereits 
sehr gut funktioniert, solange man bei funktionierendem Code nicht immer weiter nachfragt und etwas verbessern lassen will.
Angestoßen durch ein begleitendes Modul zum Thema 'Secure Software Engineering' entstand das Interesse an der Frage, inwieweit 
LLMs auch bei der Erstellung von Software-Dokumentation unterstützen können. Insbesondere stellt sich dabei die Frage, ob LLMs 
in der Lage sind, eine vollständige und aussagekräftige Dokumentation autonom zu generieren oder ob sie eher als Hilfsmittel zur 
Unterstützung menschlicher Autoren dienen können.

\section{Zielsetzung der Arbeit} \index{Zielsetzung der Arbeit} \label{Zielsetzung der Arbeit}

Das Ziel der Arbeit ist es, herauszufinden, wie gut LLM-Tools geeignet sind, Dokumente aus dem Software-Engineering-Prozess zu 
erstellen oder bei der Erstellung zu unterstützen. Der Fokus liegt dabei vor allem auf der Erstellung der Kernabschnitte der 
einzelnen Dokumente, wie zum Beispiel dem fachliche Datenmodell und den einzelnen UML-Diagrammen. Des Weiteren sollen Studenten 
einen Eindruck erhalten, wie sie die Eingaben formulieren müssen und worauf es dabei ankommt. Zusätzlich soll eine Tendenz 
abgegeben werden, welche Tools für die Erstellung besonders gut geeignet sind und welche Tools weniger geeignet erscheinen. 

\section{Struktur der Arbeit} \index{Struktur der Arbeit} \label{Struktur der Arbeit}

Grundsätzlich befasst sich die Arbeit mit drei ausgewählten LLM-Tools: ChatGPT und Google Gemini als vermutlich die 
bekanntesten Tools sowie Le Chat von Mistral AI als kleine europäische Alternative. Mit diesen Tools wird versucht, 
die Dokumente für ein ``Mensch ärgere dich nicht''-Spiel selbstständig zu erstellen. Dabei wird ausschließlich mit den 
Ausgaben der Tools gearbeitet, ohne sie zu verändern oder zu verbessern. Stattdessen wird mit Hilfe von Eingaben 
versucht, die Tools zu korrigieren. Die zu erstellenden Dokumente und die Anforderungen an das Projekt orientieren 
sich an den Vorgaben des Moduls ``Secure Software Engineering'' an der Universität der Bundeswehr München, in der 
Fakultät ``ETTI6'' am Institut für ``Software Engineering''. Die verantwortliche Professorin ist Prof. Dr. rer. nat. Andrea Baumann, 
die mit ihren Kontaktdaten im Literaturverzeichnis aufgeführt ist.

Die Arbeit ist wie folgt strukturiert: Zunächst wird eine Einführung in die Historie der drei betrachteten Tools gegeben. 
Anschließend erfolgt eine detaillierte Beschreibung der in einem Software-Engineering-Prozess typischerweise erstellten 
Dokumente, einschließlich ihrer Inhalte und ihres Zwecks. Darauf aufbauend wird untersucht, wie die eingesetzten Large 
Language Models (LLMs) bei der Erstellung dieser Dokumente unterstützen können. Dieses Kapitel formuliert zugleich die 
zentrale Fragestellung der Analyse.\\
Im folgenden Kapitel werden die Ergebnisse der praktischen Anwendung der LLMs präsentiert und vergleichend betrachtet. 
Dabei wird insbesondere auf die Qualität der generierten Inhalte sowie auf mögliche Fehler und Limitationen eingegangen.\\
Die Arbeit schließt mit einer Zusammenfassung der wichtigsten Erkenntnisse ab. Zudem werden kritisch die Schwächen der 
durchgeführten Analyse und der Arbeitsweise reflektiert sowie, in einem Ausblick, mögliche Ansatzpunkte für zukünftige Forschungsarbeiten aufgezeigt.

Die erstellten Ausgaben, welche im Kapitel 4 analysiert werden, befinden sich auf dem beigelegten USB-Stick 
in Form von PDF-Dokumenten [\autoref{BeigelegteDateien}]. In der Überschrift der Dokumente steht jeweils, 
um was es sich dabei handelt und dann in runden Klammern das Tool, welches für die Erstellung verwendet 
wurde, und eine Zahl, die wievielte Eingabe dies zu diesem Thema war. In den einzelnen Abschnitten wird als erstes der 
Prompt niedergeschrieben, der für die Erstellung der Ausgabe genutzt wurde und anschließend folgt die Ausgabe. In der Analyse wird immer auf 
das Kapitel innerhalb eines Dokuments und die Seite verwiesen. Der Verweis bezieht sich dabei immer auf 
den äquivalenten Abschnitt im beiliegenden Dokument. Als Beispiel: Wenn auf einen Inhalt im Abschnitt ``Anforderungsspezifikation'' verwiesen wird, 
ist dieser in dem Dokument ``04\_Anforderungsspezifikation.pdf'' zu finden. Zu beginn jedes Unterkapitels wird aber auch nochmal 
erwähnt, in welchem Dokument sich die Ausgaben befinden.

Im Rahmen der Arbeit wurde eine Tabelle erstellt, die als Übersicht dient. In den einzelnen Kapiteln wird diese Tabelle 
schrittweise aufgedeckt, um einen schnellen Überblick über die Ergebnisse des jeweiligen Kapitels zu liefern.

\chapter{Grundlagen} \index{Grundlagen}

TODO!!!!!!!!!

\section{LLM Tools} \index{LLM Tools} \label{LLM Tools}

Unter LLM-Tools versteht man Sprachmodelle, die auf einem Large Language Model (LLM) basieren. Ein LLM ist ein 
Deep-Learning-Algorithmus, der mit sehr großen Datensätzen trainiert wird. Diese Modelle finden häufig Anwendung im 
Bereich des Natural Language Processing (NLP), wo sie verwendet werden, um Abfragen in natürlicher Sprache zu beantworten 
oder Ergebnisse zu liefern. LLMs können neue Inhalte verstehen, zusammenfassen, generieren und vorhersagen. Durch das Training 
sammeln LLMs Milliarden von Parametern, bei denen es sich um Variablen handelt, die im Modell angepasst werden, um neue 
Inhalte abzuleiten.\\
LLMs basieren auf einem Transformer-Modell, das Eingaben in Token umwandelt und dann gleichzeitig mathematische 
Gleichungen ausführt, um Beziehungen zwischen den Token zu ermitteln. Dadurch kann der Computer Muster erkennen, 
die auch ein Mensch wahrnehmen würde, wenn ihm die gleiche Frage gestellt wird. Zudem verwenden Transformer-Modelle 
Selbstaufmerksamkeitsmechanismen, die es dem Modell ermöglichen, schneller zu lernen als herkömmliche Modelle. 
Dadurch kann das Transformer-Modell verschiedene Teile der Sequenz oder den gesamten Kontext eines Satzes berücksichtigen, 
um Vorhersagen zu generieren.

Grundsätzlich bestehen LLMs aus vier neuronalen Netzwerkschichten: der wiederkehrenden Ebene, der Einbettungsebene, 
der Feedforward-Ebene und der Aufmerksamkeitsebene. Diese Schichten arbeiten zusammen, um den Eingabetext zu verarbeiten 
und Ausgabeinhalte zu generieren.\\
Die wiederkehrende Ebene dient dazu, die Wörter des Eingabetextes der Reihe nach zu interpretieren und die Beziehungen 
zwischen den Wörtern in einem Satz zu erfassen. Die Einbettungsebene erfasst die semantische und syntaktische Bedeutung 
der Eingabe, sodass das Modell den Kontext verstehen kann. Die Feedforward-Schicht besteht aus mehreren vollständig 
verbundenen Schichten, die die Eingabeeinbettungen transformieren. Dadurch ermöglichen diese Schichten dem Modell, 
Abstraktionen auf höherer Ebene zu verstehen und somit die Absicht des Benutzers mit der Texteingabe zu erfassen. Die 
Aufmerksamkeitsebene ermöglicht es dem Sprachmodell, sich auf einzelne Teile des Eingabetextes zu konzentrieren, die für 
die Aufgabe relevant sind, und dadurch die genausten Ausgaben zu generieren.

Damit ein großes Sprachmodell Texteingaben empfangen und eine Ausgabevorhersage generieren kann, muss es zunächst 
allgemein geschult und anschließend feinabgestimmt werden, um spezifische Aufgaben ausführen zu können. Für die 
Schulung werden riesige Datenmengen im Petabyte-Bereich benötigt. Das Training verläuft mehrstufig und beginnt in 
der Regel mit einem unbeaufsichtigten Lernansatz, bei dem das Modell mit unstrukturierten und unbeschrifteten Daten 
trainiert wird, da diese in größeren Mengen verfügbar sind. In dieser Phase leitet das Modell Beziehungen zwischen 
verschiedenen Wörtern und Konzepten ab.\\
Anschließend erfolgt die Schulung und Feinabstimmung durch eine Form des selbstüberwachten Lernens. Dabei wird eine 
Datenkennzeichnung durchgeführt, durch die das Modell verschiedene Konzepte besser und genauer identifizieren kann. 
Im nächsten Schritt durchläuft das LLM den transformatorischen neuronalen Netzwerkprozess im Rahmen des Deep Learning. 
Die Transformer-Modellarchitektur ermöglicht es dem LLM, mittels eines Selbstaufmerksamkeitsmechanismus Beziehungen 
und Verbindungen zwischen Wörtern und Konzepten zu erkennen, indem bestimmten Elementen (Token) Bewertungen zugewiesen 
werden, um die Beziehungen festzulegen.

Entscheidend über die Leistungsfähigkeit und die Qualität des Sprachmodells sind die Datensätze, welche zum trainieren
benutzt werden. Im folgenden wird auf die drei unterschiedliche LLM-Tools eingeganden, welche in dieser Arbeit näher 
untersucht werden\cite{GrundlagenLLM}. 

\subsection{ChatGPT} \index{ChatGPT} \label{ChatGPT}

ChatGPT wurde von dem Unternehmen OpenAI am 30. November 2022 veröffentlicht. Aufgrund der fortschrittlichen generativen 
KI-Fähigkeiten wurde der Chatbot schnell zu einer internationalen Sensation. Die Beta-Version der Plattform war 
ursprünglich kostenlos verfügbar, und OpenAI stellte den Nutzern weiterhin eine kostenlose Basisversion der Plattform zur 
Verfügung. Im Februar 2023 kündigte OpenAI ein kostenpflichtiges Abonnement an, das den Nutzern unbegrenzten Zugriff auf 
ChatGPT und eine breitere Palette von Funktionen und Diensten bietet.
Neben begeisterter Kritik löste ChatGPT jedoch auch schnell mehrere Kontroversen aus. Es wurden fehlerbehaftete 
Antworten ausgegeben, insbesondere wenn sie aufgefordert werden, Schlussfolgerungen zu ziehen, Nuancen zu untersuchen, 
Meinungen zu analysieren oder Vorhersagen zu treffen. Außerdem äußerten sich Pädagogen auf der ganzen Welt besorgt über 
den Platz von ChatGPT im Klassenzimmer und auch an Universitäten wurde der Einsatz von ChatGPT kritisch gesehen.
Doch auch mit den Kontroversen schien ChatGPT bereit zu sein, die Art und Weise, wie Internetnutzerinhalte erstellt werden,
zu revolutionieren. Webentwickler und Ersteller von Inhalten begannen schnell, ChatGPT als Recherche- und Schreibhilfe für 
Website-Texte zu nutzen.
Im März 2023 veröffentlichte OpenAI eine neue Version, bekannt als GPT-4. Ein großer Unterschied ist, dass GPT-4 
Eingaben sowohl in text- als auch in bildbasierten Formaten akzeptiert, was dem ChatBot ermöglicht, Daten aus 
Diagrammen, Grafiken und Screenshots zu interpretieren. Ebenfalls macht GPT-4 weitaus weniger Denkfehler und sachliche 
Fehler, die in der KI als "Halluzinationen" bezeichnet werden, als das GPT-3.5-Modell, das in der ersten ChatGPT-Version 
verwendet wurde. Am Ende des Jahres war es immer fortschrittlicher geworden und ChatGPT könnte beispielsweise Bilder 
analysieren, Sprachansagen abhören und seine eigene gesprochene Antwort anbieten. Im Mai 2024 kündigte OpenAI ein 
aktualisiertes Modell namens ChatGPT-4o an, das schnellere Ergebnisse und verbesserte Text- und Audiofunktionen versprach,
einschließlich eines neuen Konversations-KI-Systems namens "Sky"\cite{GrundlagenChatGPT}.

Wenn man ChatGPT fragt, ob es einen Infotext zu sich selbst erstellen kann erhält man folgende Antwort:

\begin{figure}[H] 
    \begin{tcolorbox}[colback=gray!20, colframe=gray!20, boxrule=0pt, sharp corners] 
    \textbf{Prompt: Erstell mir einen Infotext über dich für meine Bachelorarbeit}\\ 

    ChatGPT ist ein fortschrittlicher KI-gestützter Sprachassistent, der auf der GPT-4-Architektur von OpenAI basiert. 
    Entwickelt im Rahmen der Forschung und Entwicklung im Bereich Künstliche Intelligenz, ist ChatGPT in der Lage, 
    natürliche Sprache zu verstehen und darauf zu antworten. Dies macht ihn zu einem vielseitigen Werkzeug für eine 
    Vielzahl von Anwendungen, von der Beantwortung alltäglicher Fragen bis hin zur Unterstützung bei komplexeren 
    Aufgaben wie dem Schreiben von Texten, der Generierung von Ideen und der Bereitstellung von Informationen.\\
    
    Die GPT-4-Architektur, auf der ChatGPT basiert, verwendet tiefes maschinelles Lernen, um menschliche Sprache zu 
    verarbeiten und zu erzeugen. Durch die Analyse großer Mengen an Textdaten kann das Modell Muster und Zusammenhänge 
    erkennen, die es ihm ermöglichen, kontextuell relevante und kohärente Antworten zu geben. Diese Fähigkeiten machen 
    ChatGPT zu einem nützlichen Instrument in Bereichen wie Bildung, Forschung, Kundenservice und mehr.\\
    
    Ein weiterer bemerkenswerter Aspekt von ChatGPT ist seine Fähigkeit, sich an eine Vielzahl von Themen anzupassen 
    und Informationen auf klare und präzise Weise zu vermitteln. Dies wird durch kontinuierliches Training und 
    Verbesserung des Modells erreicht, wobei sowohl technische Fortschritte als auch Nutzerfeedback berücksichtigt 
    werden.\\
    
    Zusammengefasst ist ChatGPT ein leistungsstarker und vielseitiger KI-Assistent, der durch seine fortschrittliche 
    Technologie und seine breite Anwendbarkeit einen wertvollen Beitrag in vielen Bereichen leistet.\\
\vfill
    \end{tcolorbox}
\end{figure}

Wenn man 

\subsection{Google Gemini} \index{Google Gemini} \label{Google Gemini}

Google Gemini (ehemals Bard) wurde von DeepMind entwickelt, was ein Tochterunternehmen von Google ist.  
Das Erstveröffentlichungsdatum war der 6. Dezember 2023 und basiert auf der Meena-Architektur von Google AI,  
die 2020 veröffentlicht wurde.

\subsection{Le Chat} \index{Le Chat} \label{Le Chat}

Le Chat ist ein, von dem französischem KI-Startup Mistral AI veröffentlichter Chatbot. Mistral AI startete mit  
seinen frei verwendbaren Sprachmodellen, also auf der Grundlage von Open Source, erfolgreich durch. Nun(????)  
hat Mistral AI sein bislang größtes Modell 'Mistral Large' veröffentlicht. Diesmal allerdings nicht auf der  
Basis von Open Source, sondern ausschließlich über die eigene Webseite und der KI-Infrastruktur Microsoft  
Azure. Es lassen sich allerdings API-Keys für Programmierschnittstellen erstellen, um z.B. Mistral Large über  
seinen eigenen Server laufen zu lassen und so für andere User auf der eigenen Homepage verfügbar zu machen.  
Mit Mistral Large wurde auch Le Chat veröffentlicht, welcher aktuell(?????????) kostenfrei verwendet werden kann.
Le Chat bietet derzeit noch sehr wenige Funktionen an. Es stehen Lediglich Texteingabe und -ausgabe zur Verfügung.  
Die Datengrundlage reicht nur bis 2021, weshalb es auch hier, für die Jahre 2022 bis heute, zu der Problematik der  
Halluzination kommen kann.
Grundsätzlich kann man zwischen drei Sprachmodellen auswählen: Large, Next und Small. Large bietet überlegene Denkfähigkeit,  
Next ist ein Prototyp-Modell für erhöhte Kürze und Small arbeitet schnell und kosteneffektiv. \cite{GrundlagenLeChat}

\section{Software Engineering Prozess}
Test Verweis auf Buch \cite{Sommerville10} Das ist ein Text.
\subsection{einzelne Phasen} \index{einzelne Phasen} \label{einzelne Phasen}


\cleardoublepage 
\pagenumbering{Roman}
\cleardoublepage

% Erstellen des Anhangs

\chapter{Anhang} \label{Anhang}

\section*{Kundengespräch}

\begin{prompt}[H]
    \begin{tcolorbox}[colback=gray!20, colframe=gray!20, boxrule=0pt, sharp corners] 
        **Kunde (Herr Müller):** Guten Tag, ich bin Herr Müller. Wie geht es ihnen?

        **Entwickler (Frau Schmidt):** Hallo, ich bin Frau Schmidt. Mir geht es gut. Wie geht es ihnen?
        
        **Kunde (Herr Müller):** Mir geht es auch gut. Ich möchte ein digitales "Mensch ärgere dich nicht"-Spiel entwickeln lassen und brauche Ihre Expertise dafür.
        
        **Entwickler (Frau Schmidt):** Es freut mich, dass Sie sich für unsere Dienstleistungen interessieren. Können Sie mir bitte mehr über Ihr Projekt erzählen? Welche Funktionen und Eigenschaften sollte das Spiel haben?
        
        **Kunde:** Natürlich. Zunächst möchte ich, dass das Spiel als JavaFX Anwendung auf dem Computer läuft. Die Anwendung soll eine Server-Client-Applikation, die mit Hilfe von Java realisiert wird, sein. Server und Client sollen über RMI kommunizieren. Es sollte das klassische "Mensch ärgere dich nicht"-Spielbrett und die entsprechenden Regeln enthalten. Außerdem soll es auch möglich sein, mehrere Spiele gleichzeitig zu spielen.
        
        **Entwickler:** Verstehe. Sollen die Spieler gegen den Computer, gegen andere Spieler online oder beides spielen können?
        
        **Kunde:** Beides wäre ideal. Es sollte einen Einzelspielermodus gegen KI-Gegner sowie einen Mehrspielermodus geben, in dem man online gegen Freunde oder zufällige Gegner spielen kann.
        
        **Entwickler:** Gut, dann werden wir einen Multiplayermodus integrieren. Wie viele Spieler sollen maximal teilnehmen können?
        
        **Kunde:** Wie im klassischen Spiel sollten es bis zu vier Spieler sein.
        
        **Entwickler:** Perfekt. Haben Sie Vorstellungen bezüglich des Designs? Soll das Spiel ein klassisches, traditionelles Aussehen haben oder möchten Sie ein modernes Design?
        
        **Kunde:** Ein klassisches Design mit einer ansprechenden, modernen Benutzeroberfläche wäre gut. Es sollte an das traditionelle Spiel erinnern, aber gleichzeitig frisch und ansprechend wirken.
        
        **Entwickler:** Verstanden. Gibt es besondere Funktionen oder Features, die Ihnen wichtig sind? Zum Beispiel Chat-Funktionalitäten, Punktesysteme oder spezielle Animationen?
        
        **Kunde:** Ein Chat während des Spiels wäre nett. Ein Punktesystem könnte die Motivation erhöhen, also wäre das auch eine gute Idee. Animationen sollten einfach, aber ansprechend sein - nichts Übertriebenes.
        
        **Entwickler:** Das klingt alles machbar. Wie sieht es mit dem Zeitrahmen aus? Haben Sie eine bestimmte Deadline, bis wann das Spiel fertig sein soll?
        
        **Kunde:** Ich denke, ein Zeitraum von sechs Monaten sollte machbar sein. Glauben Sie, dass das realistisch ist? Darf ich außerdem Fragen wie ihr Team aufgestellt ist?
        
        **Entwickler:** Ja, sechs Monate sind ein realistischer Zeitraum für dieses Projekt, vorausgesetzt, wir haben regelmäßige Feedback-Schleifen und Sie können uns zeitnah Ihre Anforderungen und Rückmeldungen geben. Wir sind ein Team aus 4 Personen. Ich bin die Projektleiterin, der Herr Schneider ist unser Anforderungsanalyst und fachlicher Chefdesigner, die Frau Fischer ist unsere Systemarchitektin und der Herr Becker ist der Testmanager.
        
        **Kunde:** Das klingt gut. Was wären die nächsten Schritte?
        \vfill
    \end{tcolorbox}
    \caption{Kundengespräch}
    \label{Kundengespräch1}
\end{prompt}

\begin{prompt}[H]
    \begin{tcolorbox}[colback=gray!20, colframe=gray!20, boxrule=0pt, sharp corners] 
        **Entwickler:** Wir würden zunächst einen detaillierten Projektplan erstellen, der alle Funktionen und Meilensteine festlegt. Danach beginnen wir mit der Entwicklung des Prototyps. Während der gesamten Entwicklung würden wir regelmäßige Meetings abhalten, um den Fortschritt zu besprechen und eventuelle Anpassungen vorzunehmen.
        
        **Kunde:** Das klingt nach einem guten Plan. Wie sind sie eigentlich dazu gekommen, Entwickler zu werden?
        
        **Entwickler:** Ich habe mich schon als Kind für Computerspiele interessiert und auch wie diese im Hintergrund funktionieren. Daher habe ich recht früh angefangen das Programmieren zu lernen. Das hat mir immer viel Spaß gemacht und so habe ich nach der Schule eine Ausbildung in der Branche angefangen.
        
        **Kunde:** Ah ich verstehe. Sie sind jetzt mittlerweile Projektleiterin oder?
        
        **Entwickler:** Ja, genau.
        
        **Kunde:** Fehlt ihnen die Arbeit als Programmiererin? Als Projektleitern machen Sie ja wahrscheinlich nicht mehr ganz so viel in dem Bereich oder?
        
        **Entwickler:** Ne, das stimmt. Aber man ist ja trotzdem noch damit in Berührung. Wenn mal irgendwo Schwierigkeiten bei der Implementierung auftreten, kann ich auch helfen diese zu lösen. Ein wenig fehlen tut es mir schon, aber meine jetzige Tätigkeit als Projektleiterin macht mir auch sehr viel Spaß.
        
        **Kunde:** Alles klar. Ich freue mich darauf, mit Ihnen zusammenzuarbeiten.
        
        **Entwickler:** Vielen Dank, Herr Müller. Wir freuen uns auch auf die Zusammenarbeit. Wir werden Ihnen in Kürze einen detaillierten Projektplan und ein Angebot zukommen lassen.
        
        **Kunde:** Perfekt, ich warte auf Ihre Nachricht. Einen schönen Tag noch!
        
        **Entwickler:** Ihnen auch, Herr Müller. Bis bald!
        \vfill
    \end{tcolorbox}
    \caption{Kundengespräch}
    \label{Kundengespräch2}
\end{prompt}


\section*{Tabelle der auf dem USB-Stick enthaltenen Dateien} \label{BeigelegteDateien}

\begin{longtable}{|p{5cm}|p{10cm}|}
    \hline
    \textbf{Datei} & \textbf{Beschreibung} \\
    \hline
    \endfirsthead

    01\_Besprechungsprotokolle.pdf & Enthält die erstellten Besprechungsprotokolle der drei Tools. \\
    \hline
    02\_Projekthandbuch.pdf & Enthält die erstellten Abschnitte vom Projekthandbuch der drei Tools. \\
    \hline
    03\_Gemini\_Risikoliste.xlsx & Enthält die erstellte Risiko- und Maßnahmentabelle von Gemini. \\
    \hline
    03\_GPT\_Risikoliste.xlsx & Enthält die erstellte Risiko- und Maßnahmentabelle von ChatGPT. \\
    \hline
    03\_LeChat\_Risikoliste.xlsx & Enthält die erstellte Risiko- und Maßnahmentabelle von Le Chat. \\
    \hline
    04\_Anforderungsspezifikation.pdf & Enthält die erstellten Dokumente der Anforderungsspezifikation der drei Tools. \\
    \hline
    05\_Architekturdokument.pdf & Enthält die erstellten Dokumente des Architekturdokumentes der drei Tools. \\
    \hline

\end{longtable}

\cleardoublepage

%% Abkürzungsverzeichnis

\printglossary[type=\acronymtype]
\printglossary

\listoffigures
\cleardoublepage 

%% Tabellenverzeichnis
\listoftables
\cleardoublepage 

%% Quellcodeverzeichnis
\renewcommand\lstlistlistingname{Quellcodeverzeichnis} 
\lstlistoflistings 
\renewcommand*\lstlistingname{Quellcode}
\cleardoublepage 

%% Stichwortverzeichnis, Index
\renewcommand{\indexname}{Stichwortverzeichnis}
\printindex
\cleardoublepage 

\printbibliography
\cleardoublepage 

\end{document}